\documentclass[11pt]{article}
\usepackage{amsmath,amssymb,amsthm,mathtools,bm}
\usepackage[margin=1in]{geometry}
\usepackage{hyperref}
\usepackage{enumitem}
\title{\textbf{Problem 5 Solution: Fundamental Limitations}}
\author{}
\date{}

\newtheorem{theorem}{Theorem}
\newtheorem{proposition}{Proposition}
\newtheorem{corollary}{Corollary}
\newtheorem{lemma}{Lemma}
\theoremstyle{definition}
\newtheorem{definition}{Definition}
\newtheorem{remark}{Remark}
\newtheorem{assumption}{Assumption}

\begin{document}
\maketitle

\section{Setup and scope}

We consider the translation-only model on the admissible center-of-mass region $\Omega_r \subset \mathbb{R}^3$:
\begin{equation}
 m\ddot x = F_{\mathrm{mag}}(x,u) - mg\,e_z.
 \label{eq:plant}
\end{equation}
As in Problems 1--4, $e_z=(0,0,1)^\top$ points upward, so gravity is $-mg\,e_z$.

Problem 5 asks whether every point in $\Omega_r$ can be stabilized, how spatial decay of actuator influence affects this possibility, and whether stable levitation is possible under purely static conservative field assumptions.

We organize the answer into two layers:
\begin{enumerate}[label=\textbf{Layer \arabic*:}, leftmargin=1.5em]
    \item \textbf{Abstract control layer.} Here we reason from the bounded reachable force set and the local input Jacobian, using the results of Problems 1 and 2.
    \item \textbf{Physical field layer.} Here we impose source-free static field assumptions and prove an Earnshaw-type obstruction via the maximum principle.
\end{enumerate}

\section{Layer A: abstract necessary conditions for global stabilization}

\subsection{Failure mode 1: lack of equilibrium feasibility}

For each $x\in\Omega_r$, let
\[
\mathcal{F}(x):=\{F_{\mathrm{mag}}(x,u):u\in U\}
\]
denote the bounded reachable force set.

\begin{proposition}[Global feasibility necessary condition]
If there exists a point $x^\star\in\Omega_r$ such that
\begin{equation}
 mg\,e_z \notin \mathcal{F}(x^\star),
 \label{eq:not_feasible}
\end{equation}
then $x^\star$ admits no equilibrium input, and therefore not every point in $\Omega_r$ can be stabilized.
\end{proposition}

\begin{proof}
By the equilibrium characterization from Problem 1, a point $x^\star$ is an equilibrium if and only if
\[
mg\,e_z\in\mathcal{F}(x^\star).
\]
Thus \eqref{eq:not_feasible} implies that no equilibrium exists at $x^\star$. Since local asymptotic stabilization about a point requires that point to be an equilibrium of the closed-loop system, $x^\star$ cannot be stabilized.
\end{proof}

\begin{remark}
This proposition is not deep by itself; it is the global restatement of Problem 1. Its role here is to identify the first obstruction: some points may fail before stabilizability is even considered.
\end{remark}

\subsection{Failure mode 2: equilibrium exists but is not locally stabilizable}

Suppose $x^\star\in\Omega_r$ admits an equilibrium input $u^\star$. Linearizing the translation-only dynamics about $(x^\star,u^\star)$ gives
\begin{equation}
 \delta\dot z = A\,\delta z + B\,\delta u,
 \label{eq:linsys}
\end{equation}
with state $\delta z=(\delta x,\delta v)^\top$, and
\begin{equation}
 A=
 \begin{bmatrix}
 0 & I\\
 m^{-1}A_x & 0
 \end{bmatrix},
 \qquad
 B=
 \begin{bmatrix}
 0\\
 m^{-1}B_x
 \end{bmatrix},
 \label{eq:AB}
\end{equation}
where
\[
A_x=\frac{\partial F_{\mathrm{mag}}}{\partial x}(x^\star,u^\star),
\qquad
B_x=\frac{\partial F_{\mathrm{mag}}}{\partial u}(x^\star,u^\star).
\]

\begin{proposition}[Rank deficiency as a second failure mode]
Let $x^\star$ be an equilibrium point. If there exists an eigenvalue $\lambda$ of $A$ with $\operatorname{Re}(\lambda)\ge 0$ and a nonzero left eigenvector $q^\top$ such that
\begin{equation}
 q^\top A = \lambda q^\top,
 \qquad
 q^\top B = 0,
 \label{eq:PBHfail}
\end{equation}
then the pair $(A,B)$ is not stabilizable, so $x^\star$ is not locally asymptotically stabilizable.
\end{proposition}

\begin{proof}
Condition \eqref{eq:PBHfail} is exactly the Popov--Belevitch--Hautus (PBH) failure condition for stabilizability. Therefore the linearization is not stabilizable. By the necessity of linearized stabilizability for local asymptotic stabilization of the nonlinear system, $x^\star$ is not locally asymptotically stabilizable.
\end{proof}

\begin{remark}
A useful geometric interpretation is the following: if a direction lies in the left nullspace of $B_x$, then first-order control input cannot directly affect acceleration in that direction. If the field-gradient term $A_x$ has an unstable mode there, stabilization fails even though the point is force-feasible.
\end{remark}

\section{Decay-induced loss of global force authority}

We now show that bounded actuators together with distance decay can rule out global levitation on sufficiently large domains.

\begin{assumption}[Decay model]
For each face actuator $i\in\{1,\dots,6\}$, there exists a face location $p_i\in\partial\Omega_r$, constants $C>0$ and $\alpha>0$, such that
\begin{equation}
 \|g_i(x)\| \le \frac{C}{\|x-p_i\|^\alpha}
 \qquad \text{for all } x\in\Omega_r.
 \label{eq:decay}
\end{equation}
Assume also symmetric bounded inputs $|u_i|\le u_{\max}$.
\end{assumption}

In the control-affine surrogate,
\[
F_{\mathrm{mag}}(x,u)=\sum_{i=1}^6 u_i g_i(x).
\]
Let $x_c$ denote the cube center. Then $\|x_c-p_i\|=L/2$ for every face actuator.

\begin{proposition}[Critical domain size under decay]
Under \eqref{eq:decay}, the maximum force magnitude achievable at the cube center satisfies
\begin{equation}
 \|F_{\mathrm{mag}}(x_c,u)\| \le 6u_{\max}\,\frac{C}{(L/2)^\alpha}.
 \label{eq:centerbound}
\end{equation}
Consequently, if
\begin{equation}
 6u_{\max}\,\frac{C}{(L/2)^\alpha} < mg,
 \label{eq:hoverfail}
\end{equation}
then the cube center cannot be levitated. Equivalently, defining the critical size
\begin{equation}
 L^\star := 2\left(\frac{6u_{\max}C}{mg}\right)^{1/\alpha},
 \label{eq:Lstar}
\end{equation}
any cube with side length $L>L^\star$ fails to admit hover at its center.
\end{proposition}

\begin{proof}
At the center $x_c$, the decay bound gives
\[
\|g_i(x_c)\|\le \frac{C}{(L/2)^\alpha}
\qquad \text{for each } i.
\]
Hence, for any admissible $u$,
\[
\|F_{\mathrm{mag}}(x_c,u)\|
=
\left\|\sum_{i=1}^6 u_i g_i(x_c)\right\|
\le
\sum_{i=1}^6 |u_i|\,\|g_i(x_c)\|
\le
6u_{\max}\,\frac{C}{(L/2)^\alpha},
\]
which proves \eqref{eq:centerbound}. If \eqref{eq:hoverfail} holds, then the upward force needed to cancel gravity, whose magnitude is $mg$, is outside the reachable force ball at $x_c$, so $mg\,e_z\notin\mathcal{F}(x_c)$. By Proposition 1, the center cannot be levitated. Solving \eqref{eq:hoverfail} for $L$ gives \eqref{eq:Lstar}.
\end{proof}

\begin{remark}
This proposition gives a concrete bridge between the abstract model and physics: even if a controller is perfectly designed, a sufficiently large domain can defeat hover simply because all face influences decay too much before reaching the interior.
\end{remark}

\section{Layer B: Earnshaw obstruction for static conservative fields}

We now impose genuine static field assumptions.

\begin{assumption}[Static source-free conservative model]
Fix a constant actuator input $u^\star\in U$ (no feedback, no time variation). Assume:
\begin{enumerate}[label=(\roman*), leftmargin=1.5em]
    \item the field is static and source-free in $\operatorname{int}(\Omega_r)$;
    \item there exists a scalar potential energy $\Psi:\operatorname{int}(\Omega_r)\to\mathbb{R}$ such that
    \begin{equation}
    F_{\mathrm{mag}}(x,u^\star) = -\nabla \Psi(x);
    \label{eq:consforce}
    \end{equation}
    \item the total potential
    \begin{equation}
    V(x):=\Psi(x)+mgz
    \label{eq:totalpotential}
    \end{equation}
    is harmonic in $\operatorname{int}(\Omega_r)$, i.e.
    \begin{equation}
    \Delta V = 0.
    \label{eq:harmonicV}
    \end{equation}
\end{enumerate}
\end{assumption}

\begin{remark}
Assumption \eqref{eq:harmonicV} is the mathematical form of the source-free static-field hypothesis used in Earnshaw-type arguments. Since $mgz$ is linear in $x$, it is harmonic:
\[
\Delta(mgz)=0.
\]
Therefore
\[
\Delta V = \Delta(\Psi+mgz)=\Delta\Psi,
\]
so the harmonicity of the total potential $V$ is equivalent to the harmonicity of the field-induced potential $\Psi$. The proof below uses only harmonicity of the total potential; it does not depend on the detailed actuator geometry.
\end{remark}

\begin{lemma}[No strict interior minima for harmonic potentials]
Let $V$ be harmonic on a connected open set $D\subset\mathbb{R}^3$. Then $V$ has no strict local minimum in $D$ unless it is constant on a neighborhood of that point.
\end{lemma}

\begin{proof}
This is the strong maximum principle for harmonic functions.
\end{proof}

\begin{theorem}[Earnshaw-type obstruction]
Under Assumption 2, the system admits no asymptotically stable open-loop equilibrium in $\operatorname{int}(\Omega_r)$.
\end{theorem}

\begin{proof}
An open-loop equilibrium $x^\star$ for the static input $u^\star$ must satisfy
\[
F_{\mathrm{mag}}(x^\star,u^\star)-mg\,e_z=0.
\]
Using \eqref{eq:consforce} and \eqref{eq:totalpotential}, this is equivalent to
\[
\nabla V(x^\star)=0.
\]
For the undamped conservative system \eqref{eq:plant} with fixed input $u^\star$, the total energy
\[
H(x,\dot x)=\frac12 m\|\dot x\|^2 + V(x)
\]
is conserved. Therefore asymptotic stability is impossible regardless of the shape of $V$. Moreover, even Lyapunov stability of $(x^\star,0)$ requires $x^\star$ to be a local minimum of $V$, since otherwise arbitrarily small perturbations can lower the potential and move away from $x^\star$ while preserving total energy.

But by Assumption \eqref{eq:harmonicV}, $V$ is harmonic in $\operatorname{int}(\Omega_r)$. By the strong maximum principle, a harmonic function cannot possess a strict local minimum in the interior unless it is locally constant. Hence $V$ cannot have a strict interior local minimum at $x^\star$. This rules out even Lyapunov stability, and therefore certainly rules out asymptotic stability, for any open-loop equilibrium in $\operatorname{int}(\Omega_r)$.
\end{proof}

\begin{remark}
The conclusion is the precise mathematical content of the Earnshaw obstruction in this setting: purely static, source-free, conservative fields cannot robustly trap the robot at an interior point without feedback or additional nonstatic structure.
\end{remark}

\section{Consequences}

\begin{corollary}[Necessity of feedback or non-static actuation]
If the objective is robust interior stabilization throughout a nontrivial subset of $\Omega_r$, then at least one of the following is necessary:
\begin{enumerate}[label=(\roman*), leftmargin=1.5em]
    \item active feedback control;
    \item time-varying or nonstatic fields;
    \item nonconservative or hybrid actuation (for example, onboard thrusters or other body-fixed force sources);
    \item operation outside the source-free static conservative regime covered by Theorem 1.
\end{enumerate}
\end{corollary}

\begin{proof}
The abstract results of Problems 1 and 2 show that global stabilization requires both equilibrium feasibility and local stabilizability. Proposition 3 shows that bounded decaying authority can destroy feasibility on large domains. Theorem 1 shows that even when an equilibrium exists, purely static conservative source-free fields cannot make it asymptotically stable in open loop. Therefore robust stabilization requires leaving that regime, for example through feedback, time variation, or hybrid actuation.
\end{proof}

\section{Final modeling-fidelity remark}

\begin{remark}[What the abstract surrogate can bypass]
The abstract surrogate
\[
F_{\mathrm{mag}}(x,u)=G(x)u
\]
is useful for control design because it exposes force authority and stabilizability in a transparent way. However, it can represent force fields with properties that physically realizable magnetostatic systems do not possess, such as arbitrary full-rank force maps everywhere in the interior or force directions incompatible with source-free static field structure. In particular, the surrogate can satisfy rank-$3$ conditions at every point even though a true static electromagnetic system is still subject to the Earnshaw obstruction and Maxwell constraints. Thus the surrogate is a control-design model, not a guarantee of physical realizability.
\end{remark}

\section{Final summary}

The formal solution to Problem 5 is:
\begin{enumerate}[label=\textbf{(\arabic*)}, leftmargin=1.5em]
    \item Global stabilization can fail because some points are not even force-feasible: if $mg\,e_z\notin\mathcal{F}(x^\star)$, no equilibrium exists there.
    \item Even when an equilibrium exists, local stabilization can fail if the linearized pair $(A,B)$ is not stabilizable; the PBH test gives the precise obstruction.
    \item Under explicit decay bounds, there is a critical domain size $L^\star$ above which even the cube center cannot be levitated.
    \item Under static source-free conservative field assumptions with harmonic total potential, no asymptotically stable open-loop interior equilibrium can exist; this is the Earnshaw-type obstruction.
    \item Therefore global robust stabilization requires feedback, time variation, hybrid actuation, or departure from the static source-free conservative regime.
\end{enumerate}

\end{document}
